%%═══════════════════════════════════════════════════════════════════
%% Lecture 04 — Jacobi Convergence, Multidimensional Roots, and Minimization
%% 77315 — Computational Physics
%% Date: 26/04/2026
%%═══════════════════════════════════════════════════════════════════

\section{Lecture 04 — Jacobi Convergence, Multidimensional Roots, and
  Minimization}\label{sec:lec04}

\begin{center}
\begin{tabular}{ll}
  \textbf{Date:}\quad 26/04/2026 & \\
\end{tabular}
\end{center}
\hrule\vspace{1.5em}

%%──────────────────────────────────────────────────────────────────
\subsection{Jacobi Method: Convergence Rate and Complexity}

\textbf{Convergence Rate per Cycle.} Define a \textbf{cycle} as one complete
pass through all $\frac{n(n-1)}{2}$ off-diagonal element pairs. Let
$S^{(k)} = \sum_{i \neq j} A_{ij}^2$ after $k$ cycles. Applying the
per-rotation bound repeatedly over a full cycle:
\begin{equation}
S^{(1)} \leq S^{(0)} \cdot \left(1 - \frac{2}{n(n-1)}\right)^{n(n-1)/2}
\approx S^{(0)} \cdot e^{-1}
\end{equation}
Each cycle reduces the off-diagonal weight by $\approx e^{-1} \approx 0.37$.
In practice, 10 cycles is sufficient for nearly all applications.

\textbf{Computational Complexity.} Each rotation modifies 2 rows and 2
columns of $n$ elements, requiring $\sim 12$ multiplications per zero
created. A full cycle has $\frac{n(n-1)}{2}$ rotations; for 10 cycles the
total is $\approx 60n^3$, confirming $\bigO{n^3}$ scaling.

\textbf{Tracking Eigenvectors.} Maintain the accumulated rotation matrix
$T = T_1 T_2 \cdots T_k$ (initialised as $T = I$). After each plane rotation
with $\cos\phi = c$, $\sin\phi = s$ on the $(p,q)$ plane, update:
\begin{align}
T_{r,p} &\leftarrow c\, T_{r,p} - s\, T_{r,q} \\
T_{r,q} &\leftarrow s\, T_{r,p} + c\, T_{r,q}
\end{align}
At convergence the columns of $T$ are the eigenvectors $\varphi_i$.

\begin{tcolorbox}[colback=blue!5!white,colframe=blue!75!black,
  title=Optimisation: Row-Maximum Vector ($\bigO{n^4} \to \bigO{n^3}$)]
\textbf{Naive bottleneck}: scanning the full upper-triangular part to find
$\max_{i<j}|A_{ij}|$ costs $\bigO{n^2}$ per rotation, giving total
$\bigO{n^4}$.

\textbf{Fix}: maintain a vector $\mathtt{j\_max}[i]$ storing the column index
of the largest off-diagonal element in row $i$. Finding the global max then
costs $\bigO{n}$. After each rotation (which modifies rows and columns $p$
and $q$), only rows $p$, $q$, and rows that had their max in $\{p,q\}$ need
re-scanning — costing $\bigO{n}$ expected total per rotation. Net result:
$\bigO{n^3}$ total.
\end{tcolorbox}

%%──────────────────────────────────────────────────────────────────
\subsection{Finding the Root of a Function in Multiple Dimensions}

\noindent\textbf{Physical intuition:} In 1D, bracketing (bisection) works
because a sign change guarantees a root. In $N$ dimensions, the zero-set of
each component $F_i(\mathbf{x}) = 0$ is a hypersurface; roots are
intersections of $N$ such hypersurfaces. There is no bounding box whose
corner values can certify a root lies inside — we must use local iterative
methods that depend sensitively on the starting point.

\textbf{Multidimensional Newton-Raphson.} Construct the flat tangent
hyperplane at the current point using the \textbf{Jacobian matrix}
$J_{ij} = \frac{\partial F_i}{\partial x_j}$. The Newton step $\delta$
satisfies:
\begin{equation}
\boxed{J\,\delta = -F(\mathbf{x})}
\end{equation}
This is solved by Gaussian elimination (Lecture~\ref{sec:lec02}). The new
estimate is $\mathbf{x}_{\text{new}} = \mathbf{x} + \delta$, and the process
repeats until $|F(\mathbf{x})|$ is acceptably small. Convergence is
\emph{quadratic} near a root (correct digits double each iteration), but a
poor starting point can diverge.

%%──────────────────────────────────────────────────────────────────
\subsection{Line Search and Backtracking}\label{subsec:linesearch}

\begin{tcolorbox}[colback=yellow!5!white,colframe=yellow!60!black,
  title=The Blind Hiker Analogy]
A blindfolded hiker finds the downhill direction via Newton's method but
taking a massive step might overshoot the valley bottom and ascend the
opposite face. Solution: define the scalar merit function
$f(\mathbf{x}) = \frac{1}{2}F(\mathbf{x}) \cdot F(\mathbf{x})$ and take a
step $\lambda p$. If $f(\mathbf{x}_{\text{new}}) > f(\mathbf{x})$, backtrack
by reducing $\lambda$ using parabolic/cubic curve fitting. Every step must
strictly lower $f$.
\end{tcolorbox}

\textbf{Descent Guarantee.} With the Newton direction $p = -J^{-1}F$:
\begin{equation}
\grad f \cdot p = F^T J \cdot (-J^{-1}F) = -F \cdot F < 0
\end{equation}
so $p$ is always a descent direction. We accept $\lambda$ once the
\textbf{Armijo condition} is satisfied:
\begin{equation}
f(\mathbf{x} + \lambda p) \leq f(\mathbf{x}) + \varepsilon\,\grad f(\mathbf{x})\cdot(\lambda p),
\qquad \varepsilon \approx 10^{-4}
\end{equation}

\textbf{Backtracking Steps.}
\begin{tcolorbox}[colback=blue!5!white,colframe=blue!75!black,
  title=Parabolic and Cubic Polynomial Fitting]
\textbf{First backtrack — parabolic fit.} From $g(0)$, $g'(0) < 0$, and
$g(1)$, the parabola minimum is at:
\begin{equation}
\lambda^* = -\frac{g'(0)}{2\bigl[g(1) - g(0) - g'(0)\bigr]}
\end{equation}
Clamped to $[0.1,\, 0.5]$ to prevent steps that are negligibly small or
dangerously large.

\textbf{Subsequent backtracks — cubic fit.} With four known values
$g(0)$, $g'(0)$, $g(\lambda_1)$, $g(\lambda_2)$, solve for the cubic minimum:
\begin{equation}
\frac{1}{\lambda_1 - \lambda_2}
\begin{pmatrix} \lambda_1^{-2} & -\lambda_2^{-2} \\ -\lambda_1^{-3} &
\lambda_2^{-3} \end{pmatrix}
\begin{pmatrix} g(\lambda_1) - g(0) - g'(0)\lambda_1 \\
g(\lambda_2) - g(0) - g'(0)\lambda_2 \end{pmatrix}
= \begin{pmatrix} a \\ b \end{pmatrix}
\end{equation}
Minimum at $\lambda^* = \bigl(-b + \sqrt{b^2 - 3a\,g'(0)}\bigr)/(3a)$,
again clamped.
\end{tcolorbox}

%%──────────────────────────────────────────────────────────────────
\subsection{Finding a Minimum in One Dimension --- Golden Section Search}
To bracket the minimum of a 1D function without derivatives, maintain three
points $a < b < c$. Evaluate a new point $x$ inside the larger interval to
shrink the bracket. The mathematically optimal placement guarantees equal
fractional reduction regardless of outcome, forcing the intervals to obey
the \textbf{Golden Ratio}:
\begin{equation}
w = \frac{3 - \sqrt{5}}{2} \approx 0.382
\end{equation}
By always choosing $x$ exactly $38.2\%$ into the larger segment, the bracket
shrinks optimally at every step.

%%──────────────────────────────────────────────────────────────────
\subsection{Minimization in a Multidimensional Space}
When finding a multidimensional minimum, gradient descent leads to a flat
spot ($\grad f = 0$). To confirm it is a minimum (not a saddle point), the
landscape must curve upward in \emph{every} direction: the \textbf{Hessian
Matrix} $H_{ij} = \frac{\partial^2 f}{\partial x_i \partial x_j}$ must be
\textbf{Positive Definite}.

\noindent\textbf{Efficient check:} Apply Gaussian elimination to $H$. If and
only if every diagonal pivot $B_{ii} > 0$ during elimination, the matrix is
positive definite, confirming a true minimum. (Computing all eigenvalues
would cost $\bigO{n^3}$ with a large prefactor; this check is cheaper and
equally conclusive.)

%%──────────────────────────────────────────────────────────────────
\subsection*{Recommended Reading}
\begin{itemize}
    \item \textit{Numerical Recipes}, W.H. Press et al.\ --- Ch.\ 11.1
    (Jacobi Transformations), Ch.\ 9.7 (Globally Convergent Methods for
    Nonlinear Systems), Ch.\ 10.1 (Golden Section Search).
    \item \textit{Computational Physics}, Mark Newman --- Ch.\ 6 (Solution of
    Linear and Nonlinear Equations).
\end{itemize}
